\documentclass[11pt]{article}
\usepackage[margin=2.1cm]{geometry}
\usepackage{amsmath,amssymb,amsthm,mathrsfs}
\usepackage{longtable,array}
\usepackage{url}
\usepackage[hidelinks]{hyperref}
\numberwithin{equation}{section}
\theoremstyle{plain}
\newtheorem{theorem}{Theorem}[section]
\newtheorem{proposition}[theorem]{Proposition}
\newtheorem{lemma}[theorem]{Lemma}
\newtheorem{corollary}[theorem]{Corollary}
\theoremstyle{definition}
\newtheorem{definition}[theorem]{Definition}
\newtheorem{example}[theorem]{Example}
\newtheorem{remark}[theorem]{Remark}
\newtheorem{hypothesis}[theorem]{Hypothesis}
\allowdisplaybreaks
\emergencystretch=3em
\newcommand{\Ap}{A|_{\perp}}
\newcommand{\AV}{A|_V}
\newcommand{\Ob}{\mathrm{Ob}}
\newcommand{\ICT}{I^{CT}[\varepsilon]}
\newcommand{\lc}{\mathrm{lc}}
\newcommand{\dslash}{{\not d}}

\title{The action-level one-loop equivariant quantum master equation\\ along every generator of $\mathfrak{psl}(4|4)$}
\author{Step 2, strict Costello route, action side (theory)}
\date{2026-10-30}

\begin{document}
\maketitle

\begin{abstract}
For $\mathcal N=4$ super Yang--Mills theory in the formalism of [C07] on Euclidean $\mathbb R^4$, with the heat-kernel cutoff $\varepsilon$ on every edge and the transfer endpoint $T$, we set up and analyse the one-loop equivariant quantum master equation of Definition c8-quant at the level of the action, arity by arity, along every generator: the dilatation $\Delta$, the supercharges $\xi\in S_+\otimes W$ and $\xi\in S_-\otimes W^*$, the special supercharges $\eta\in\Sigma^*$, the special conformal fields $K_b$, and, as regressions, the translations $P_b$, the rotations $\mathfrak{so}(V)$ and the R-symmetry $\mathfrak{so}(W^\perp_{\mathbb R})$. The obstruction $\Ob$ of Definition c8-ob is written component by component as the $\varepsilon\to0$ limit of the defect graphs. We prove: the one-loop counterterm cochain is $\tfrac94$ units times $\varepsilon^{-1}\int(A,A)$ plus $\log\varepsilon$ times an $s$-closed cochain (no pole in arity one or two); every component $\Ob^{(1)}(x)$ and $\Ob^{(2)}(x,y)$ is the sum of the $\varepsilon^0$ coefficients of the irreducible defect cycles with one special edge, the vertex $S^{(1)}(x)$ (and $S^{(1)}(y)$ or $S^{(2)}(x,y)$) placed on the cycle in the listed ways, all other classes cancelling or converging to trees; a second, independent evaluation applies $Q$ to the renormalised arity-one and arity-two amplitudes, which the conjugation of $S^{(1)}(x)$ by $e^{\hbar\partial_P}$ (terminating at second order along every generator, the cubic term of $S^{(1)}(\xi)$ included) organises into variation kernels, composite lines and endpoint terms. The census of the sectors that weight, ghost number, the centres of $\mathrm{Spin}(4)\times SU(4)$ and colour allow is given for all components. The Wess--Zumino identities are derived in every arity. The arity-one equation is solvable exactly when a residual class in $H^0(Q+\{I,-\})$ vanishes; along $\Sigma$ the two-leg cohomology is zero, so the two-leg equation is solvable with the two free directions $S^{(1)}(\xi)$ restricted to the $V$-arrows and to the $W^\perp_{\mathbb R}$-arrows. The cubic part of the arity-one equation fixes the differences of the sector coefficients exactly when the sector parts of $S^{(1)}(x)$ fail to extend to symmetries of the interacting classical theory. The arity-two residuals form a two-cocycle of $\mathfrak{psl}(4|4)$ with values in $H^{-1}(Q+\{I,-\})$; under the hypothesis that the classical symmetries of local degree $-1$ are the superconformal ones, the arity-two equations fix the uniform coefficients $t_x$ of $S^{(1)}(x)$ in $\hbar F_1^{(1)}(x)$ up to the inner direction generated by $S^{(1)}(\Delta)$ and the outer direction $G_B$; in particular the coefficient $\theta$ of Theorem c9-kappa is fixed by the two-leg part of the arity-two equation along $(\Delta,P_b)$, $2u\,\theta_d$ equal to the coefficient of the class of $S^{(1)}(P_b)_d$ in the arity-two obstruction. The one-loop corrections to $[Q,Q]$, $\{Q,S\}$, $[K,P]$, $[D,Q]$ are the components of this cocycle. Closed forms (simplex integrals with the special edge, generating functions, momentum injection for the position) and the specification of two independent engine routes with their gates complete the note; the scaling identity, $K=P^\dagger$, the twist and the value $\theta=-\tfrac13$ enter as cross-checks only.
\end{abstract}

\tableofcontents

\section{Data, names and inputs}
\label{sec:at-data}

The source of truth is the manuscript, Chapters 1--11; labels such as Definition c8-ob refer to it. The unit is
\[
u:=\frac{\hbar\operatorname{tr}_{\mathfrak g}(\mathrm{ad}_X\mathrm{ad}_Y)}{16\pi^2(X,Y)}=-\frac{g^2N}{32\pi^2}\quad(\mathfrak g=\mathfrak{su}(N),\ (X,Y)=-\tfrac{4\hbar}{g^2}\operatorname{tr}(XY),\ \text{Proposition c8-twocomp}).
\]
Edges carry $P(\varepsilon,T)=\int_\varepsilon^T(Q^{GF}\otimes1)K_t\,dt$ (Definition c8-PT); $T=\infty$ along every generator except $K_b$, and along the pairs $(K_b,P_{b'})$ and $(\eta,\eta')$ with $\Gamma(\eta,\eta')|_V\neq0$, where $T$ is finite (Proposition c8-Tsafe, Definition c8-K). No dimensional regularisation and no infrared regulator other than $T$ is used anywhere.

\begin{definition}[Generators and their data]
\label{def:at-gens}
For a homogeneous generator $x$ of $\mathfrak{psl}(4|4)$ let $\lambda_x$ be its $\mathrm{ad}_\Delta$-eigenvalue ($[\Delta,x]=\lambda_xx$, Remark c2-dilatation), $p_x$ the polynomial degree of its Killing data in the position, and $|x|$ its parity. The functional $S^{(1)}(x)$ (Definition c2-S1, \eqref{eq:at-S1list} below) has ghost degree $-1$ and Weyl weight $4-\tfrac12\lambda_x$, the position counted at weight $-1$.
\[
\begin{array}{l|c|c|c|l|l}
x & \lambda_x & p_x & |x| & \text{terms of }S^{(1)}(x) & \text{endpoint}\\
\hline
\Delta\ (v=2E) & 0 & 1 & 0 & \text{quadratic: }(v.A,A^*),\ -(v.\lambda,\lambda^*),\ -(v.c,c^*) & T=\infty\\
P_b & -2 & 0 & 0 & \text{quadratic} & T=\infty\\
\mathfrak{so}(V)\ni M_\omega\ (v=\omega x) & 0 & 1 & 0 & \text{quadratic} & T=\infty\\
\mathfrak{so}(W^\perp_{\mathbb R})\ni R_r & 0 & 0 & 0 & \text{quadratic, on }\Ap,\lambda\text{ only} & T=\infty\\
\xi\in S_+\otimes W\ \text{or}\ S_-\otimes W^* & -1 & 0 & 1 & T_1,\dots,T_4\ \text{quadratic},\ T_5\ \text{cubic} & T=\infty\\
\eta\in S_+\otimes W^*\ \text{or}\ S_-\otimes W & 1 & 1 & 1 & \textstyle\sum_\mu x^\mu S^{(1)}(\rho(e_\mu)\eta)+\int s_0 & T=\infty\\
K_b\ (v=(x,x)b-2(b,x)x) & 2 & 2 & 0 & \text{quadratic} & T<\infty
\end{array}
\]
Here, with $T_1,\dots,T_5$ the terms of Table c11-terms and $s_0=-(\rho(\Ap)\eta,\lambda^*)$ ((c11-Seta) of Chapter 11),
\begin{equation}
\begin{aligned}
&S^{(1)}(x)=\int\mathrm{dvol}\,\big((v.A,A^*)-(v.\lambda,\lambda^*)-(v.c,c^*)\big)\quad(x\ \text{even}),\\
&S^{(1)}(\xi)=T_1+T_2+T_3+T_4+T_5,\qquad T_1=-\int(\Gamma(\xi,\lambda)|_V,A^*|_V),\ T_2=-\int(\Gamma(\xi,\lambda)|_\perp,A^*|_\perp),\\
&T_3=\tfrac12\int(\rho(d\AV)\xi,\lambda^*),\ T_4=\tfrac12\int(\rho(d\Ap)\xi,\lambda^*),\ T_5=\tfrac12\int(\rho(\tfrac12[A\wedge A])\xi,\lambda^*).
\end{aligned}
\label{eq:at-S1list}
\end{equation}
The leg names used throughout: $\lambda_+=\lambda|_{S_+\otimes W}$, $\lambda_-=\lambda|_{S_-\otimes W^*}$, $\lambda^*_+=\lambda^*|_{S_+\otimes W^*}$, $\lambda^*_-=\lambda^*|_{S_-\otimes W}$, $\xi_+\in S_+\otimes W$, $\xi_-\in S_-\otimes W^*$, $\eta_+\in S_+\otimes W^*$, $\eta_-\in S_-\otimes W$.
\end{definition}

\begin{definition}[Components, sector parts, coefficients]
\label{def:at-names}
\leavevmode
\begin{itemize}
\item (i) A cochain $f$ of $C^\bullet(\mathfrak{psl}(4|4),\mathscr O_l(\mathscr E))$ has components $f^{(k)}(x_1,\dots,x_k)$ (Definition c2-cochains), normalised as in Lemma c2-components: the arity-one component of $d_{CE}f+\tfrac12\{f,f\}$ at $f=S_{\mathfrak{psl}(4|4)}$ is $\{S,S^{(1)}(x)\}$ and its arity-two component is $\{S,S^{(2)}(x,y)\}+S^{(1)}([x,y])+(-1)^{|x|}\{S^{(1)}(x),S^{(1)}(y)\}$. $\Lambda^{(0)}:=Q+\{I,-\}=\{S,-\}$ on local functionals (Theorem c8-L (iv)). $\Ob^{(k)}$ and $F_1^{(k)}$ are the components of the obstruction and of the finite counterterm $\hbar F_1$ in this normalisation.
\item (ii) Sector parts. For even $x$ the sector parts of $S^{(1)}(x)$ are
\[
S^{(1)}(x)_\perp=\int(v.\Ap,A^*|_\perp),\qquad S^{(1)}(x)_\lambda=-\int(v.\lambda,\lambda^*),\qquad S^{(1)}(x)_V=\int(v.\AV,A^*|_V)-\int(v.c,c^*),
\]
with sum $S^{(1)}(x)$; each is $Q$-closed, $v.$ commuting with $Q$ (Propositions c7-cov, c7-dil). For $\xi\in\Sigma$, $S^V_2(\xi):=T_1+T_3$ and $S^\perp_2(\xi):=T_2+T_4$ (Definition c11-terms); for $\eta\in\Sigma^*$, $S^V(\eta):=\sum_\mu\int x^\mu S^V_2(\rho(e_\mu)\eta)$ and $S^\perp(\eta):=\sum_\mu\int x^\mu S^\perp_2(\rho(e_\mu)\eta)+\int s_0$, whose sum is the quadratic part of $S^{(1)}(\eta)$. The index $d$ runs over the sectors $\{\perp,\lambda,V\}$ (even $x$) or $\{V,\perp\}$ (odd $x$).
\item (iii) The quadratic part of $\hbar F_1^{(1)}(x)$ contains $u\sum_d\theta^x_d$ times the sector parts; $\theta^\Delta_\perp,\theta^\Delta_\lambda,\theta^\Delta_V$ are $\theta,\theta_\lambda,\theta_V$ of Definition c11-free and $\theta^{P_b}_d$ are $\theta_P,\theta_{P,\lambda},\theta_{P,V}$; along $\xi$, $\theta^\xi_V,\theta^\xi_\perp$ are those of Proposition qs-freexi. When all $\theta^x_d$ are equal their common value is the uniform coefficient $t_x$, the coefficient of $u\,S^{(1)}(x)$ (completed by its cubic part).
\item (iv) $\ICT$ is the counterterm cochain $I^{CT}_{\mathfrak{psl}(4|4)}[\varepsilon]$ of Definition c8-quant; its coefficient of $\log\varepsilon$ is the cochain $\lc$, and $\lc^{(k)}$ its components.
\end{itemize}
\end{definition}

\begin{remark}[Physics notation]
\label{rem:at-phys}
$S_+$ carries undotted and $S_-$ dotted indices; $W$ and $W^*$ are the $\mathbf 4$ and $\bar{\mathbf 4}$ of $SU(4)_R$ (upper and lower index $i$), $W^\perp_{\mathbb R}$ the vector $\mathbf 6$ of $SO(6)_R$, realised by the antisymmetric $A_{ij}=A_I\Sigma^I_{ij}$ (Definition c8-twocomp); $\rho$ is Clifford multiplication, in components the ten-dimensional chiral $\Gamma^M$, in four dimensions $\sigma^m,\bar\sigma^m$, and on the Dirac column $\rho(v)=-i\gamma_S^\mu v_\mu=\gamma_W^\mu v_\mu$ (Srednicki, Weinberg). The gaugino $\lambda$ lies in $\Pi\Sigma\otimes\mathfrak g$, $\Pi$ the parity shift, $\Sigma=S_+\otimes W\oplus S_-\otimes W^*$ the ten-dimensional chiral spinor $\mathbf{16}$, $\Sigma^*$ the $\mathbf{16}'$ (lower ten-dimensional index). Functional degrees (ghost numbers of the coordinate functionals): $c$: $+1$, $A$, $\lambda$: $0$, $A^*$, $\lambda^*$: $-1$, $c^*$: $-2$; Weyl weights $0,1,\tfrac32,3,\tfrac52,4$. For $\mathfrak g=\mathfrak{su}(N)$, $A_m=-igA^a_mT^a$, $\lambda=-\tfrac{ig}{\sqrt2}\lambda^aT^a$, $\operatorname{tr}T^aT^b=\tfrac12\delta^{ab}$ (Proposition c8-twocomp).
\end{remark}

\begin{remark}[Inputs used as established]
\label{rem:at-inputs}
From the manuscript: the arity-zero results (Theorems c4-quadratic, c4-log, c4-obstruction, c4-finite, c9-ob3, c9-F3, c9-wz), the two-leg data along $\Delta$ of Theorem c9-quadratic and Remark c9-Pi as regression targets (their derivation in Chapter 9 uses the scaling identity; here they are to be recomputed through the defect graphs), the commutators of Lemma c11-lines and Proposition qs-comm, the simplex theorems c10-simplex, c11-simplex, c11-oddsimplex, and the endpoint statements of Chapter 8. From the earlier engines: their sign rules and normalisations (Section \ref{sec:at-engines}).
\end{remark}
